% Part: proof-theory
% Chapter: normalization
% Section: introduction

\documentclass[../../../include/open-logic-section]{subfiles}

\begin{document}

\olfileid{pt}{nor}{int}

\olsection{导论}

如果你已经推导出了一个条件句~$!A \lif !B$，你也许可以借助~\Elim{\lif} 在推导~$!B$ 的过程中使用它。
这当然也需要一个~$!A$ 的推导~$\delta_1$。
你关于~$!A \lif !B$ 的推导~$\delta_1$，十有八九以~\Intro{\lif} 结束，该规则将一个从开放假设~$!A$ 出发推导~$!B$ 的推导~$\delta_2$，转化为~$!A \lif !B$ 的推导。
如果确实如此，你最终的推导就会呈现如下形式：
\[
\AxiomC{$\Discharge{!A}{x}$}
\RightLabel{$\delta_2$}
\DeduceC{$!B$}
\DischargeRule{\Intro{\lif}}{x}
\UnaryInfC{$!A \lif !B$}
\AxiomC{}
\RightLabel{$\delta_2$}
\DeduceC{$!A$}
\RightLabel{\Elim{\lif}}
\BinaryInfC{$!B$}
\DisplayProof
\]
尽管你的策略是高效的，因为你复用了已有的东西（一个~$!A \lif !B$ 的推导），但你的推导本身却并非如此。
先引入~$!A \lif !B$ 只是为了随即消去它，这似乎是在绕路。
应当存在一种更直接的方法来推导~$!B$。

事实上，总是存在这样的方法。
像上面这样的“绕路”——即引入规则之后紧跟着消去规则——总是可以从自然演绎的推导中消除。
这被称为\emph{规范化}，其结果是一个\emph{正规}推导（或称处于\emph{正规形}的推导）。

与相继式演算中的切割消去定理类似，这个结果的证明有些复杂，需要我们考虑大量情形。
对经典系统~\Log{N1c} 的证明比直觉主义系统~\Log{N1i} 更难。
正如切割消去表明在相继式演算中使用~\Cut{} 规则并不能推导出更多东西一样，规范化表明使用绕路也不会比避免绕路证明出更多的东西。
还有其他相似之处：对于同一结果，正规推导可能比包含绕路的最短推导长得多；正如在相继式演算中，带切割的推导可能比最短的无切割推导短得多。
\emph{子公式性质}——即你总可以仅用结果的子公式来推导该结果——在自然演绎中由规范化得出，正如在相继式演算中由切割消去得出一样。

\end{document}